\documentclass{book}
\usepackage[backend=biber,natbib=true,style=authoryear]{biblatex}
\addbibresource{/home/nguyen/1_NQBH/math/bib.bib}
\usepackage[utf8]{inputenc}
\usepackage{float}
\usepackage{graphicx}
\usepackage[colorlinks=true,linkcolor=blue,urlcolor=red,citecolor=magenta]{hyperref}
\usepackage{amsmath,amssymb,amsthm}
\allowdisplaybreaks
\numberwithin{equation}{section}
\newtheorem{assumption}{Assumption}[section]
\newtheorem{lemma}{Lemma}[section]
\newtheorem{corollary}{Corollary}[section]
\newtheorem{definition}{Definition}[section]
\newtheorem{proposition}{Proposition}[section]
\newtheorem{theorem}{Theorem}[section]
\newtheorem{notation}{Notation}[section]
\newtheorem{remark}{Remark}[section]
\newtheorem{example}{Example}[section]
\newtheorem{ques}{Question}[section]
\newtheorem{problem}{Problem}[section]
\newtheorem{conjecture}{Conjecture}[section]
\usepackage[left=0.5in,right=0.5in,top=1.5cm,bottom=1.5cm]{geometry}
\usepackage{fancyhdr}
\pagestyle{fancy}
\fancyhf{}
\addtolength{\headheight}{0pt}% obsolete
\lhead{\scriptsize \chaptername~\thechapter}
\rhead{\scriptsize \nouppercase{\leftmark}} %\nouppercase !
\renewcommand{\chaptermark}[1]{\markboth{#1}{}}
\cfoot{\thepage}
\def\labelitemii{$\circ$}

\title{Giftedness}
\author{Collector: Nguyen Quan Ba Hong}
\date{\today}

\begin{document}
\maketitle
\setcounter{secnumdepth}{6}
\setcounter{tocdepth}{6}
\tableofcontents

%------------------------------------------------------------------------------%

\part{Wikipedia's}

\chapter{\href{https://en.wikipedia.org/wiki/Eidetic_memory}{Wikipedia/Eidetic Memory}}

\textit{Eidetic memory} (eye-DET-ik; more commonly called \textit{photographic memory}) is the ability to recall an image from \href{https://en.wikipedia.org/wiki/Memory}{memory} with high precision for a brief period after seeing it only once,[1] and without using a \href{https://en.wikipedia.org/wiki/Mnemonic_device}{mnemonic device}.[\href{http://www.britannica.com/EBchecked/topic/180955/eidetic-image}{Eidetic image $|$ psychology}, Encyclop\ae dia Britannica online]

Although the terms \textit{eidetic memory} and \textit{photographic memory} are popularly used interchangeably,[1] they are also distinguished, with \textit{eidetic memory} referring to the ability to see an object for a few minutes after it is no longer present[2][``\href{http://plato.stanford.edu/entries/mental-imagery/quasi-perceptual.html}{Mental Imagery $>$ Other Quasi-Perceptual Phenomena (Stanford Encyclopedia of Philosophy)}''. plato.stanford.edu. Retrieved 2016-04-30.] and \textit{photographic memory} referring to the ability to recall pages of text or numbers, or similar, in great detail.
\begin{itemize}
    \item Anthony Simola (2015). \href{https://books.google.com/books?id=p4ZMBwAAQBAJ&pg=PT117}{\textit{The Roving Mind: A Modern Approach to Cognitive Enhancement}}. ST Press. p. 117. ISBN 978-0692409053. Retrieved May 10, 2016.
    \item ``\href{http://www.slate.com/articles/health_and_science/science/2006/04/kaavya_syndrome.single.html}{No one has a photographic memory}''. Slate Magazine. Apr 2006.
\end{itemize}
When the concepts are distinguished, eidetic memory is reported to occur in a small number of children and generally not found in adults,[2][Andrew Hudmon (2009). \href{https://books.google.com/books?id=GKdUyxrqL6kC&pg=PA52}{Learning and Memory}. Infobase Publishing. p. 52. ISBN 978-1438119571. Retrieved May 10, 2016.] while true photographic memory has never been demonstrated to exist.
\begin{itemize}
    \item ``\href{http://www.scientificamerican.com/article/i-developed-what-appears-to-be-a-ph/}{Does Photographic Memory Exist?}''. Scientific American.
\end{itemize}
%
The word \textit{eidetic} comes from the Greek word  (pronounced [\^e:dos], eidos) ``visible form''.[``Eidetic''. American Heritage Dictionary, 4th ed. 2000. Retrieved 2007-12-12.]

\section{Eidetic or photographic memory}
The terms \textit{eidetic memory} and \textit{photographic memory} are commonly used interchangeably,[1] but they are also distinguished.[4][5]

Scholar Annette Kujawski Taylor stated,
\begin{quotation}
    ``In eidetic memory, a person has an almost faithful mental image snapshot or photograph of an event in their memory.
    
    However, eidetic memory is not limited to visual aspects of memory and includes auditory memories as well as various sensory aspects across a range of stimuli associated with a visual image.''
\end{quotation}
[Annette Kujawski Taylor (2013). \textit{Encyclopedia of Human Memory} [3 volumes]. ABC-CLIO. p. 1099. ISBN 978-1440800269. Retrieved May 10, 2016.]

Author Andrew Hudmon commented:
\begin{quotation}
    ``Examples of people with a photographic-like memory are rare.
    
    Eidetic imagery is the ability to remember an image in so much detail, clarity, and accuracy that it is as though the image were still being perceived.
    
    It is not perfect, as it is subject to distortions and additions (like episodic memory), and vocalization interferes with the memory.''
\end{quotation}
[Andrew Hudmon (2009). \textit{Learning and Memory}. Infobase Publishing. p. 52. ISBN 978-1438119571. Retrieved May 10, 2016.]

%
``Eidetikers'', as those who possess this ability are called, report a vivid \href{https://en.wikipedia.org/wiki/Afterimage}{afterimage} that lingers in the visual field with their eyes appearing to scan across the image as it is described.
\begin{itemize}
    \item Searleman, Alan; Herrmann, Douglas J. (1994). \textit{Memory from a Broader Perspective}. McGraw-Hill. p. 313. ISBN 9780070283879.
    \item ``\href{https://www.psychologytoday.com/articles/200603/the-truth-about-photographic-memory}{The Truth About Photographic Memory}''. Psychology Today. Retrieved 2016-04-30.
\end{itemize}
Contrary to ordinary mental imagery, eidetic images are externally projected, experienced as ``out there'' rather than in the mind.

Vividness and stability of the image begins to fade within minutes after the removal of the visual stimulus.[3]

\href{https://en.wikipedia.org/wiki/Scott_Lilienfeld}{Lilienfeld} et al. stated, ``\textit{People with eidetic memory can supposedly hold a visual image in their mind with such clarity that they can describe it perfectly or almost perfectly}$\ldots$, just as we can describe the details of a painting immediately in front of us with near perfect accuracy.''[Scott Lilienfeld, Steven Jay Lynn, Laura Namy, Nancy Woolf, Graham Jamieson, Anthony Marks, Virginia Slaughter (2014). \textit{Psychology: From Inquiry to Understanding}. Pearson Higher Education. p. 353. ISBN 978-1486016402. Retrieved May 10, 2016.]

%
By contrast, photographic memory may be defined as the ability to recall pages of text, numbers, or similar, in great detail, without the visualization that comes with eidetic memory.[4]

It may be described as the ability to briefly look at a page of information and then recite it perfectly from memory.

This type of ability has never been proven to exist and is considered popular myth.[5][7]

\section{Prevalence}
Eidetic memory is typically found only in young children, as it is virtually nonexistent in adults.[5][6] 

Hudmon stated, ``Children possess far more capacity for eidetic imagery than adults, suggesting that a developmental change (such as acquiring language skills) may disrupt the potential for eidetic imagery.''[6]

Eidetic memory has been found in 2 to 10 percent of children aged 6 to 12.

It has been hypothesized that language acquisition and verbal skills allow older children to think more abstractly and thus rely less on visual memory systems.

Extensive research has failed to demonstrate consistent correlations between the presence of eidetic imagery and any cognitive, intellectual, neurological or emotional measure.[``\href{http://journals.cambridge.org/action/displayAbstract?fromPage=online&aid=7176716}{Behavioral and Brain Sciences - Abstract - 20 years of haunting eidetic imagery: where's the ghost?}''.]

%
A few adults have had phenomenal memories (not necessarily of images), but their abilities are also unconnected with their intelligence levels and tend to be highly specialized.

In extreme cases, like those of \href{https://en.wikipedia.org/wiki/Solomon_Shereshevsky}{Solomon Shereshevsky} and \href{https://en.wikipedia.org/wiki/Kim_Peek}{Kim Peek}, memory skills can reportedly hinder social skills.[Barber, Nigel (December 22, 2010). ``\href{http://www.psychologytoday.com/blog/the-human-beast/201012/remembering-everything}{Remembering everything? Memory searchers suffer from amnesia!}''. Psychology Today. Sussex. Retrieved Jul 10, 2013.]

Shereshevsky was a trained \href{https://en.wikipedia.org/wiki/Mnemonist}{mnemonist}, not an eidetic memoriser, and there are no studies that confirm whether Kim Peek had true eidetic memory.

%
According to \href{https://en.wikipedia.org/wiki/Herman_Goldstine}{Herman Goldstine}, the mathematician \href{https://en.wikipedia.org/wiki/John_von_Neumann}{John von Neumann} was able to recall from memory every book he had ever read.[Goldstine, Herman (1980). \textit{The Computer from Pascal to von Neumann}. Princeton University Press. p. 167. ISBN 0-691-02367-0.]

\section{Skepticism}
Skepticism about the existence of eidetic memory was fueled around 1970 by Charles Stromeyer, who studied his future wife, Elizabeth, who claimed that she could recall poetry written in a foreign language that she did not understand years after she had 1st seen the poem.

She also could seemingly recall random dot patterns with such fidelity as to combine 2 patterns from memory into a stereoscopic image.
\begin{itemize}
    \item Stromeyer, C. F.; Psotka, J. (1970). ``The detailed texture of eidetic images''. \textit{Nature}. 225 (5230): 346--49. doi:10.1038/225346a0. PMID 5411116. S2CID 4161578.
    \item Thomas, N.J.T. (2010). Other Quasi-Perceptual Phenomena. In \textit{The Stanford Encyclopedia of Philosophy}.
\end{itemize}
She remains the only person documented to have passed such a test.

However, the methods used in the testing procedures could be considered questionable (especially given the extraordinary nature of the claims being made),[Blakemore, C., Braddick, O., \& Gregory, R.L. (1970). \textit{Detailed Texture of Eidetic Images: A Discussion}. Nature, 226, 1267--1268.] as is the fact that the researcher married his subject.

Additionally, that the tests have never been repeated (Elizabeth has consistently refused to repeat them)[Foer, Joshua (Apr 27, 2006). ``\textit{Kaavya Syndrome: The accused Harvard plagiarist doesn't have a photographic memory. No one does}''. Slate. Retrieved Dec 16, 2012.] raises further concerns for journalist \href{https://en.wikipedia.org/wiki/Joshua_Foer}{Joshua Foer} who pursued the case in a 2006 article in \href{https://en.wikipedia.org/wiki/Slate_(magazine)}{\textit{Slate}} magazine concentrating on cases of unconscious plagiarism, expanding the discussion in \href{https://en.wikipedia.org/wiki/Moonwalking_with_Einstein}{\textit{Moonwalking with Einstein}} to assert that, of the people rigorously scientifically tested, no one claiming to have long-term eidetic memory had this ability proven.
\begin{itemize}
    \item Foer, Joshua (April 27, 2006). ``\textit{Kaavya Syndrome: The accused Harvard plagiarist doesn't have a photographic memory. No one does}''. Slate. Retrieved Dec 16, 2012.
    \item Stromeyer III, Charles (1970). ``\textit{Adult Eidetiker}'' (PDF). Psychology Today: 76--80.
\end{itemize}
%
American \href{https://en.wikipedia.org/wiki/Cognitive_scientist}{cognitive scientist} \href{https://en.wikipedia.org/wiki/Marvin_Minsky}{Marvin Minsky}, in his book \href{https://en.wikipedia.org/wiki/Society_of_Mind}{\textit{The Society of Mind}} (1988), considered reports of photographic memory to be an ``unfounded myth,''[Minsky, Marvin (1998). Society of Mind. Simon \& Schuster. p. 153. ISBN 978-0-671-65713-0. ``$\ldots$ \textit{we often hear about people with `photographic memories' that enable them to quickly memorize all the fine details of a complicated picture or a page of text in a few seconds. So far as I can tell, all of these tales are unfounded myths, and only professional magicians or charlatans can produce such demonstrations}.''] and that there is no scientific consensus regarding the nature, the proper definition, or even the very existence of eidetic imagery, even in children.[3]

%
Lilienfeld et al. stated: ``Some psychologists believe that eidetic memory reflects an unusually long persistence of the iconic image in some lucky people''.

They added: ``More recent evidence raises questions about whether any memories are truly photographic (Rothen, Meier \& Ward, 2012).

\fbox{Eidetikers' memories are clearly remarkable, but they are rarely perfect.}

Their memories often contain \fbox{minor errors, including information that was not present in the original visual stimulus}.

So even eidetic memory often appears to be reconstructive''.
\begin{itemize}
    \item Scott Lilienfeld, Steven Jay Lynn, Laura Namy, Nancy Woolf, Graham Jamieson, Anthony Marks, Virginia Slaughter (2014). \textit{Psychology: From Inquiry to Understanding}. Pearson Higher Education. p. 353. ISBN 978-1486016402. Retrieved May 10, 2016.
    \item Note: \href{https://en.wikipedia.org/wiki/Reconstructive_memory}{Reconstructive memory} is a theory of memory recall.
\end{itemize}
%
\href{https://en.wikipedia.org/wiki/Skeptical_movement}{Scientific skeptic} author \href{https://en.wikipedia.org/wiki/Brian_Dunning_(author)}{Brian Dunning} reviewed the literature on the subject of both eidetic and photographic memory in 2016 and concluded that there is ``a lack of compelling evidence that eidetic memory exists at all among healthy adults, and no evidence that photographic memory exists.

But there's a common theme running through many of these research papers, and that's that the difference between ordinary memory and exceptional memory appears to be 1 of degree.''[Dunning, Brian. ``Skeptoid \#452: Photographic Memory''. Skeptoid. Retrieved 30 Oct 2016.]

\section{Trained mnemonists}
To constitute photographic or eidetic memory, the visual recall must persist without the use of mnemonics, expert talent, or other cognitive strategies.

Various cases have been reported that rely on such skills and are erroneously attributed to photographic memory.[Joshua Foer - \textit{Moonwalking with Einstein: The Art and Science of Remembering Everything}, 2011]

%
An example of extraordinary memory abilities being ascribed to eidetic memory comes from the popular interpretations of \href{https://en.wikipedia.org/wiki/Adriaan_de_Groot}{Adriaan de Groot}'s classic experiments into the ability of chess \href{https://en.wikipedia.org/wiki/Grandmaster_(chess)}{grandmasters} to memorize complex positions of chess pieces on a chess board.

Initially, it was found that these experts could recall surprising amounts of information, far more than nonexperts, suggesting eidetic skills.

However, when the experts were presented with arrangements of chess pieces that could never occur in a game, their recall was no better than the nonexperts, suggesting that they had developed an ability to organize certain types of information, rather than possessing innate eidetic ability.

%
Individuals identified as having a condition known as \href{https://en.wikipedia.org/wiki/Hyperthymesia}{hyperthymesia} are able to remember very intricate details of their own personal lives, but the ability seems not to extend to other, non-autobiographical information.[``\href{http://www.psychologytoday.com/blog/quirks-memory/201301/people-extraordinary-autobiographical-memory}{People with Extraordinary Autobiographical Memory}''. Psychology Today.]

They may have vivid recollections such as who they were with, what they were wearing, and how they were feeling on a specific date many years in the past.

Patients under study, such as \href{https://en.wikipedia.org/wiki/Jill_Price}{Jill Price}, show brain scans that resemble those with \href{https://en.wikipedia.org/wiki/Obsessive-compulsive_disorder}{obsessive-compulsive disorder}.

In fact, Price's unusual autobiographical memory has been attributed as a byproduct of compulsively making journal and diary entries.

Hyperthymestic patients may additionally have depression stemming from the inability to forget unpleasant memories and experiences from the past.[``\href{https://www.npr.org/blogs/health/2013/12/18/255285479/when-memories-never-fade-the-past-can-poison-the-present}{When Memories Never Fade, The Past Can Poison The Present}''. NPR.org. Dec 27, 2013.]

It is a misconception that hyperthymesia suggests any eidetic ability.

%
Each year at the \href{https://en.wikipedia.org/wiki/World_Memory_Championships}{World Memory Championships}, the world's best memorizers compete for prizes.

None of the world's best competitive memorizers has a photographic memory, and no one with claimed eidetic or photographic memory has ever won the championship.[Joshua Foer - \textit{Moonwalking with Einstein: The Art and Science of Remembering Everything}, 2011]

\section{Notable claims}
Main article: \href{https://en.wikipedia.org/wiki/List_of_people_claimed_to_possess_an_eidetic_memory}{List of people claimed to possess an eidetic memory}.

%
There are a number of individuals whose extraordinary memory has been labeled ``eidetic'', but it is not established conclusively whether they use \href{https://en.wikipedia.org/wiki/Mnemonic}{mnemonics} and other, non-eidetic memory-enhancement.

`Nadia', who began drawing realistically at the age of 3 is autistic and has been closely studied.

During her childhood she produced highly precocious, repetitive drawings from memory, remarkable for being in perspective (which children tend not to achieve until at least adolescence) at the age of 3, which showed different perspectives on an image she was looking at.

E.g., when at the age of 3 she was obsessed with horses after seeing a horse in a story book she generated numbers of images of what a horse should look like in any posture.

She could draw other animals, objects and parts of human bodies accurately, but represented human faces as jumbled forms.
\begin{itemize}
    \item Selfe, Lorna; Selfe, Lorna (1977), \textit{Nadia : a case of extraordinary drawing ability in an autistic child}, Academic Press, ISBN 978-0-12-635750-9
    \item New Scientist, 1 Dec 1977, Vol. 76, No. 1080 p.577 ISSN 0262-4079
    \item Selfe, Lorna; ProQuest (Firm) (2012), \textit{Nadia Revisited : A Longitudinal Study of an Autistic Savant}, Taylor and Francis, ISBN 978-0-203-82576-1
\end{itemize}
Others have not been thoroughly tested, though savant \href{https://en.wikipedia.org/wiki/Stephen_Wiltshire}{Stephen Wiltshire}
\begin{itemize}
    \item Daniel A. Weiskopf (2017) An ideal disorder? Autism as a psychiatric kind, \textit{Philosophical Explorations}, 20:2, 175-190, doi:10.1080/13869795.2017.1312500
    \item Rebecca Chamberlain, I. C. McManus, Howard Riley, Qona Rankin \& Nicola Brunswick (2013) \textit{Local processing enhancements associated with superior observational drawing are due to enhanced perceptual functioning, not weak central coherence}, The Quarterly Journal of Experimental Psychology, 66:7, 1448-1466, doi:10.1080/17470218.2012.750678
    \item Gillian J. Furniss (2008) Celebrating the Artmaking of Children with Autism, \textit{Art Education}, 61:5, 8--12, doi:10.1080/00043125.2008.11518990
\end{itemize}
can look at a subject once and then produce, often before an audience, an accurate and detailed drawing of it, and has drawn entire cities from memory, based on single, brief helicopter rides; his 6-metre drawing of 305 square miles of New York City is based on a single 20-minute helicopter ride.
\begin{itemize}
    \item ``\href{http://news.bbc.co.uk/1/hi/health/1211299.stm}{Unlocking the brain's potential}''. BBC News. Mar 10, 2001. Retrieved Nov 8, 2007.
    \item ``Like a Skyline Is Etched in His Head''. The New York Times. Oct 27, 2009. Retrieved Feb 23, 2013.
\end{itemize}
Another less thoroughly investigated instance is the art of \href{https://en.wikipedia.org/wiki/Winnie_Bamara}{Winnie Bamara}, an Australian indigenous artist of the 1950s.[35]

\section{See also}
\begin{itemize}
    \item \href{https://en.wikipedia.org/wiki/Ayumu_(chimpanzee)}{Ayumu} - a \href{https://en.wikipedia.org/wiki/Common_chimpanzee}{chimpanzee} whose performance in short-term memory tests is higher than university students
    \item \href{https://en.wikipedia.org/wiki/Exceptional_memory}{Exceptional memory} - scientific background to the research into exceptional memory
    \item \href{https://en.wikipedia.org/wiki/Visual_memory}{Visual memory} - a form of memory which preserves some characteristics of our senses pertaining to visual experience
    \item \href{https://en.wikipedia.org/wiki/Funes_the_Memorious}{\textit{Funes the Memorious}} - a short story by \href{https://en.wikipedia.org/wiki/Jorge_Luis_Borges}{Jorge Luis Borges} discussing the \fbox{consequences of eidetic memory}
    \item \href{https://en.wikipedia.org/wiki/Omniscience}{Omniscience} - particularly in Buddhism where adepts gain capacity to know ``the 3 times'' (past, present and future)
    \item \href{https://en.wikipedia.org/wiki/Synaptic_plasticity}{Synaptic plasticity} - ability of the strength of a synapse to change\hfill$\square$
\end{itemize}

%------------------------------------------------------------------------------%

\printbibliography[heading=bibintoc]
\end{document}